% 原书第 73 页（PDF 第 86 页）
\BookPage{73}{86}
\BookHeading{盛清扬州「名医」李炳的医疗生涯及其历史记忆}{盛清扬州「名医」李炳的医疗生涯及其历史记忆\\[3pt]—兼论清代医生医名的获取与流传}{余新忠*}
十八世纪的扬州，是一个令人艳羡的繁华富庶之地，不仅富商巨贵云集，文坛学界钜子辈出，而且还汇聚了一大批颇有成就的艺人和戏子。然而令人多少有些意外的是，在如此繁盛的的时代，如此富庶的地区，在医学界，却并未像其他领域那样，涌现出声名显赫的一流人才。在为数不多的几位近代以来仍被记忆的盛清时代医家中，李炳当属其中最具影响的医生了。李炳（1729-1805年，又名李钧），字振声，号西垣，江苏仪征人。从后人的视角来看，他在扬州医界的地位，从其死后不久刊刻出版的《重修扬州府志》中其传记的篇幅之长不难看出一斑²。不过考诸史迹，他在有生之年，却似乎只是位并不「行于世」，甚至贫苦潦倒的医者。那么这种不协调以及其身后有关他的历史记忆是怎样形成的呢？对此的探讨，显然有助于我们更好地理解清代医生医名的获取与流传。

虽然李炳之名及其相关的事迹，在目下医史研究者有关温病学的论着中和历史研究者有关焦循或扬州地方史的探讨中，多有提及或简略介绍，不过专门的探讨，除了医史界一篇专门探讨其论着《辨疫琐言》的医学思想的短文外³，几付阙如。故此，笔者拟依据不算多的相关记载，尽可能对李炳的医疗人生及其身后的历史记忆做一钩沉，以期对我们更好地认识和了解盛清扬州的医疗状况以传统社会医生医名的获取与流传有所裨益。

\vfill
\begin{BookNotes}
\BookNote{*}{作者任职于南开大学中国社会史研究中心。}
\BookNote{1.}{关于盛情扬州的繁盛，可参阅安东篱着，李霞译，《说扬州：1550-1850年的一座中国城市》，中华书局2007年版，第三、四部分。}
\BookNote{2.}{其传记载嘉庆《重修扬州府志》卷五十四《人物·艺术》（中国方志丛刊，华中地方·第154号，台北：成文出版社有限公司1974年影印嘉庆十五年刊本），第4219-4223页。在该书收录的医生传记中，其篇幅仅略次于主要生活在17世纪的名医郑重光（1638-1716年，第4206-4211页），远超出其他医者。}
\BookNote{3.}{苏颖、鞠煜洁，《李炳〈辨疫琐言〉医学思想探析》，《山西中医学院学报》，2007年第3期，第12-13页。}
\end{BookNotes}

% 原书第 74 页（PDF 第 87 页）
\BookPage{74}{87}
\BookSubheading{一、主要相关史籍}

根据前述《重修扬州府志》的记载，李炳着有《金匮要略注》二十二卷、《治疫琐言》（即《辨疫琐言》）一卷和《西垣诊籍》二卷等三部著作，但除《辨疫琐言》外，其余两种均未见有刊本，至今也见有相关著述和各大图书馆的书目著录，应该已经亡佚4。《辨疫琐言》目前所见最早的刊本为年裘庆元于1936年编辑出版的「珍本医书集成」本。是为李炳流传后世的唯一著述。除此之外，至今可见的时人有关李炳的记载，亦相当有限，而且大多是著名学者焦循留下的。现对其中主要的载籍略作介绍。

就管见所及，现存文献中最早载有李炳资讯的文献，为初刊于乾隆末年(1795）的《扬州画舫录》以「李钧」之名记录了他为治愈江春族人之事。并述及他撰有《金匮要略注》一书6。

稍后，焦循在其仅存的嘉庆元年的日记中，记载了他因为儿病返回扬州，请李钧诊治之事，该日记除记载此事外，还记录了一些有关他治病神效的事迹。被刘建臻视为是焦循后来撰着的《李翁医记》的初稿。九年后，即嘉庆十年（1805)，李炳不辞劳苦，竭尽心力治愈焦循儿子与儿媳的重病，令焦循甚为感动，遂仿太史公「仓公列传」之例，撰成了《李翁医记》（两卷），成为存留至今有关李炳医疗史迹最重要的资料。该著于嘉庆十三年（1808）在阮元的资助下，会同焦循的《北湖

\vfill
\begin{BookNotes}
\BookNote{4.}{著名学者刘毓崧在咸丰十年(1860)为李炳弟子汪近垣所着的《金匮要略阐义》写序言时，就说《金匮要略注》「今已无传」。（刘毓崧，《通义堂文集》卷十一《汪近垣系那是金匮要略阐义序》，「续修四库全书」1546册影印民国刘氏刻求恕斋丛书本[光绪十四年初刊]，上海古籍出版社2002年版，第513页。}
\BookNote{5.}{冯尔康先生曾在《清代仪征人才的兴起及原因》一文提到他的三部著作（《辨疫琐言》一卷、《西垣诊籍》2卷、《金匮要略注》22卷)，载《扬州研究一江都陈轶群先生百年冥诞纪念论文集》，台北联经出版公司1996年，第558页。}
\BookNote{6.}{李斗，《扬州画舫录》卷十二《桥东录》，中华书局1960年[1795]版，第279页。}
\BookNote{7.}{焦循，《理堂日记》，姚氏咫进斋抄本，国家图书馆善本库藏。}
\BookNote{8.}{刘建臻，《焦循著述新证》，社会科学文献出版社2005年版，第73-74页。其实该日记所载只是《李翁医记》事迹中的少数几个，要说是《李翁医记》初稿，似属勉强，至多只能说是该书部分内容的初稿。}
\BookNote{9.}{焦循，《与汪损之书》、《李翁医记》，载刘建臻点校，《焦循诗文集》(下)，广陵书社2009年版，第445-446、450-457页。}
\end{BookNotes}

% 原书第 75 页（PDF 第 88 页）
\BookPage{75}{88}
\noindent 小志》一起刊刻印行。在焦循完成该着后不久，李炳就发病并很快不治，这令焦循甚为悲伤，又为其写下了颇注入感情的墓志铭1°。这又成了今人了解李炳生平事迹最重要的史料。后来，焦循在所着的笔记《忆书》中，再次谈到李炳的医术11。

再五年，嘉庆十五年，焦循参与编纂的《重修扬州府志》出版，李炳被收入该书的「人物·艺术」之中，篇幅颇长，记录的内容除了摘录焦循为其所写的墓志铭和《李翁医记》的内容外，也记录了四个焦循以往著述中未曾提及的医案12。此后，于道光三十年梓行的《重修仪征县志》也在「人物·艺术」中收入了李炳的传记，不过其内容基本是抄录《扬州府志》，并未见有新的讯息。

在李炳去世两年后，焦循得危疾，幸由炳门人汪近垣治愈，这让焦循之子廷琥记起了李炳数年前曾治愈己病情形及其其治病之理与清初侨寓扬州的名医郑重光的医案中论述暗合，遂写下了《郑素圃医案序》一文13。咸丰十年（1860），李炳弟子汪近垣之子少垣谋刊刻其父之着《金匮要略阐义》，求序于同邑著名学者刘毓崧，刘氏便依据焦氏父子的有关李炳的记载以及近垣之书的内容，撰成《汪近垣先生金匮要略阐义序》一文14。

时人留下的有关李炳（李钧）的基本资讯主要就这些，下面就主要依据这些有限的史料对李炳身前的医疗生涯即身后被记忆的历史做一探究。

\BookSubheading{二、医疗人生钩沉}

李炳生于雍正七年（1729），卒于嘉庆十年（1805），享年77岁（虚岁，下同）。这在史料中记载明确，并无疑义。他籍贯仪征，但从现存的相关史料来看，

\vfill
\begin{BookNotes}
\BookNote{10.}{焦循，《名医李君墓志铭》，载刘建臻点校，《焦循诗文集》(上)，第393-395页。}
\BookNote{11.}{焦循，《忆书》卷四、五，丛书集成初编本，中华书局1985年版，第28、38页。该书焦氏生前并未刊行，直到光绪十年才有赵之谦刊出。(参阅来新夏，《清人笔记随录》，中华书局2004年版，第301-302页）}
\BookNote{12.}{嘉庆《重修扬州府志》卷五十四《人物·艺术》，第4219-4223页。}
\BookNote{13.}{焦廷琥，《蜜梅花馆文录·郑素圃医案序》，丛书集成初编本，中华书局1985年版[1824年]，第16页。}
\BookNote{14.}{刘毓崧，《通义堂文集》卷十一《汪近垣系那是金匮要略阐义序》，第512-513页。}
\end{BookNotes}

% 原书第 76 页（PDF 第 89 页）
\BookPage{76}{89}
\noindent 似乎看不出其活动与仪征有何关联，焦循曾在其《忆书》中称他，「不善逢迎，故不行于城，常往来于淮，而往邵伯镇者最久。」15而嘉庆府志的传记则云:「晚年多寓邵伯镇、瓜洲、北湖」16。从中大概可以看出至少其晚年活动范围大概是府城周边的乡镇，而非在仪征。因此很有可能他生活的轨迹主要在府城一带。另外，关于他的名字，现在的通行的说法都为李炳，而且我们今天所见史料，比如他的《辨疫琐言》的题名，焦循所写的墓志铭以及府志中的传记，也均称其为「李炳」，不过如前所述，较早成书于乾隆末年的《扬州画舫录》则称「李钧」，而焦循保留下来的次年的日记，亦以「李钧」相称，而其字、号则都一致。如果说《扬州画舫录》偶一提到他，尚有将其名弄错的可能，那么凭焦循与他的深厚交情，断不可能将记错其名。而且「钧」的本意与「音乐」有关，与其字相配程度明显超过「炳」。所以，依笔者的推测，其本名即为李钧，炳之名，乃是嘉庆元年以后，即在他晚年因为某种缘由而新改的。

下面即按时代顺序将其的习医治病的人生经历做一钩沉。

李炳终其一生，乃一介布衣，未获任何功名，但其有一定的文化素养，有著作问世，或可称之为「儒医」。不过与大多数的「儒医」因举业不成而学医有所不同，李炳自幼就拜师学医17，并「幼习三世之书」18，似乎自小就以医生为自己的人生目标。从《李翁医记》记录的故事中可以看出，其师尚算开明，一次他偷偷治愈了一个人其师未能治愈病人的病后，师「治酒」询问缘由，了解情况后，「曰：『子悟在出我上，何可在弟子列。』急趋之悬牌。」19。不过其早年的学医经历似乎不算顺利，属于大器晚成的一类。根据焦循所撰的墓志铭的说法，他幼习医书，「苦不能得其蕴，乃学《易》，十年而有得」20。他究竟跟随其师学了多少年后才获自己行医的资格，我们不得而知，但从以上的说法中，我们至少可以看出，其成

\vfill
\begin{BookNotes}
\BookNote{15.}{焦循，《忆书》卷五，第38页。}
\BookNote{16.}{嘉庆《重修扬州府志》卷五十四《人物·艺术》，第4220页。}
\BookNote{17.}{焦循，《李翁医记》卷下，载刘建臻点校，《焦循诗文集》（下），第455页。}
\BookNote{18.}{焦循，《名医李君墓志铭》，载刘建臻点校，《焦循诗文集》（上），第393页。}
\BookNote{19.}{焦循，《李翁医记》卷下，载刘建臻点校，《焦循诗文集》(下)，第455页；焦循，《理堂日记》，嘉庆元年七月初十日。}
\BookNote{20.}{焦循，《名医李君墓志铭》，载刘建臻点校，《焦循诗文集》（上)，第395页。}
\end{BookNotes}

% 原书第 77 页（PDF 第 90 页）
\BookPage{77}{90}
\noindent 才独立行医已是成年以后之事了，因为一者其学成在十年以上，二者从他能偷偷去看其师未能治愈的病人，而其师又治酒询问缘由等情形来看，他当时应当是成年人了。

在他壮年时，应该就是三四十岁之时，李炳曾因岁末避人而至苏州，苏州有病人病咳，「寐则咳，醒则已」，曾经吴医张亮葵治之不愈，李炳治以川椒，第二天咳就停止了21。另据《理堂日记》记载，他在成功治愈了江春之弟的疾病后，有洪姓商人慕名「以五百年金聘李之楚，居楚者二年」22。这就印证了其墓志铭中有关他「尝往来吴越、荆楚之间」23的说法。

乾隆四十四年（1779)，对李炳来说，是极为重要的一年，在这一年的夏天，他结识了对他日后产生重要影响的焦循。当年夏天，焦循之父病臂痛，被医者认为将成偏枯，后经其同学祖父的介绍，向李炳求医，炳开了「黑豆半升，蚕砂二两，为末服之」药方，成功地治愈了循父的病痛。焦循因此记住了李炳，但李炳当时并未对年轻的焦循留下太多的印象24。

十一年后，即乾隆五十五年（1790），李炳再次成功地治愈了焦循呕血之病。但当时，焦循仍未对焦循的医技持有特别的信任25。不过在随后的交往中，焦循对李炳的信任可谓与日俱增，据其嘉庆元年（1796）的日记记载，当年七月，他因「儿病」从浙江返回扬州，到扬州后，便请李炳诊视，并在随后的日记中述其行医神妙之事26。可见焦此时已对李的医技颇为认同。第二年，焦循怀孕中的妻子，「忽呕逆不已，每呕必厥，日十数度」。经族中医者治疗无效，又请李炳诊治，炳不畏物议，坚持使用桂枝、干姜和黄连等药，最终治愈了焦妻之病，令焦循颇为感慨。第二年，焦妻产女后第二天，呕逆之旧疾复发，遂再请李炳的诊治，李认为症

\vfill
\begin{BookNotes}
\BookNote{21.}{焦循，《李翁医记》卷上，载刘建臻点校，《焦循诗文集》（下)，第450页。}
\BookNote{22.}{焦循，《理堂日记》，嘉庆元年七月初十日。}
\BookNote{23.}{焦循，《名医李君墓志铭》，载刘建臻点校，《焦循诗文集》(上)，第393页。}
\BookNote{24.}{焦循，《李翁医记》卷下，载刘建臻点校，《焦循诗文集》（下)，第457页。}
\BookNote{25.}{焦循，《李翁医记》卷上，载刘建臻点校，《焦循诗文集》(下)，第450页。}
\BookNote{26.}{焦循，《理堂日记》，嘉庆元年七月初十日。}
\end{BookNotes}

% 原书第 78 页（PDF 第 91 页）
\BookPage{78}{91}
\noindent 虽类似，病实不同，乃治以甘草、芍药和阿胶等药，应手而愈27。这些使得焦循对李之医技的信任程度不断加深。

乾隆六十年（1795），书生周金声因为患头痛病，久治不愈，而求医于李炳，被他治愈后，遂跟他学医。周生学医似乎颇有成绩，到李炳去世时，不过十年，周已经「以医噪于时」。但第二年，就因旧疾病复发而去世28。除了周金声外，从文献中还可以看到李炳另有一位更好的弟子汪近垣。根据刘毓崧的说法，李炳去世时(1805)，汪年仅18岁，故其跟随李炳学艺，当不会早于周。咸丰十年，其子少垣以其著作请序于刘毓崧时，应已经过世。刘是这样叙述汪的事迹的:「当是时（嘉庆十二），先生年甫弱冠，而医术已精，其后擅重名者，垂五十年，齿弥高而学弥粹。」29可见汪近垣约过世于1857年（咸丰七）。

嘉庆五年(1800)，李炳完成当今唯一存世的著作《辨疫琐言》，交请焦循批评，焦随即让其子廷琥「抄录一本，藏于家塾」30。此后，焦循对该着做了编辑修缮，誊抄清本交给李炳，自己则留有底稿31。此前，至少乾隆末年之前，李炳已经撰成《金匮要略注》32，但这部较早的作品显然未能付梓，在他去世后不久，焦循在给汪损之的信中云，他也只见过该书一半，据李炳生前说，该书存在汪损之处。尽管焦循努力谋划这两部著作的出版事宜，33但其愿望并没有实现。

嘉庆九年（1804），七十六高龄的李炳喜得贵子，而这似乎是他唯一的子

\vfill
\begin{BookNotes}
\BookNote{27.}{焦循，《李翁医记》卷上，载刘建臻点校，《焦循诗文集》（下)，第450-451页。}
\BookNote{28.}{焦循，《李翁医记》卷下，载刘建臻点校，《焦循诗文集》（下)，第455页；焦循，《忆书》卷四，第28页。}
\BookNote{29.}{刘毓崧，《通义堂文集》卷十一《汪近垣系那是金匮要略阐义序》，第512页。}
\BookNote{30.}{李炳，《辨疫琐言》，「焦廷琥附识」，载裘庆元辑，《珍本医书集成·内科类》，上海科学技术出版社1986年版，第12页。}
\BookNote{31.}{焦循，《与汪损之书》，载刘建臻点校，《焦循诗文集》(下)，第445页。}
\BookNote{32.}{李斗的《扬州画舫录》曾在记录李钧的条目中称:「所着有《金匮要略注》，多发前人所未发。(李斗，《扬州画舫录》卷十二《桥东录》，第279页)这说明当时已经完成该着，并具有一定的名声。}
\BookNote{33.}{焦循，《与汪损之书》，载刘建臻点校，《焦循诗文集》(下)，第445页。}
\end{BookNotes}

% 原书第 79 页（PDF 第 92 页）
\BookPage{79}{92}
\noindent 嗣34。

嘉庆十年（1805)六月，焦循的幼孙因族人误药而死。一个月后，他的独子廷琥又不幸身染重病，当李炳得悉情况后，感觉此子乃焦氏血胤所系，甚是着急，不顾自己年事已高，竭尽全力予以救治。对此，焦循记叙道：

\begin{BookQuote}
时闰六月二十五日，翁清晨至曰：君之孙已为医误，此子所关甚重，然病情隐曲，今终夜思之，前此非所治也，当由心阴伤而心阳上越，姑试以甘温。署甘草、大枣等令服，未服而身亦有疹大如戎豆，色且紫，他医议用快斑发疹之剂。翁又至曰：脉弦微而不渴，何敢用凉药？且未有疹出而躁若此者。是时躁甚，坐卧行立皆不宁。翁曰：试以前药服之。服己而躁定。翁曰：未也。侯之良久，果又躁，且呼手足不仁，脐下亦不仁，渐及于胃腕间。翁曰：急矣，吾今日必愈此疾。乃去急治药，促煎之。跣足袒衣，自调其水火，诊脉凡七八次，药熟又诊脉，久之自持药令服。曰：是矣，服之必愈。…
\end{BookQuote}

并最终治愈了廷琥的重病，令焦循甚是感佩，称:「余于此始恍然于忌之、谤之者真为庸医，而翁之医真能神也。」35至此，焦循对李炳的医技已是敬佩之至。

从上面的描述和他老来得子的情况看，李炳当时虽然年事已高，但身体应该还相当硬朗。然而，令人意想不到的事，就在他成功治愈焦循儿子即儿媳之病后不久，自己也突然染疾，并很快不治，于七月中离开了人世。「卒之日，予家人儿女咸哀泣。湖中农人有泣于路者。」36

\vfill
\begin{BookNotes}
\BookNote{34.}{焦循，《名医李君墓志铭》，载刘建臻点校，《焦循诗文集》(上)，第394页。在现存的相关史料中，几乎没有提到他家庭的情况的。墓志铭一般都会谈到传主的子嗣情况，但李炳的传文中只是说，「君卒年七十七，遗孤止二岁」。李炳死后，焦循对其遗孤颇为在意，死后不久，就与其朋友汪损之商量他的抚养问题。（焦循，《与汪损之书》，载刘建臻点校，《焦循诗文集》（下)，第445页）嘉庆十六年（1811年）又在给汪近垣的信中谈到「世弟李生」的学费问题，说他曾每年贴李生的读书之费，但再贴补明年的读书费，「力有不足」，故请汪近垣处理。（焦循，《与汪近垣》，载刘建臻点校，《焦循诗文集》（下)，第677页）虽然我们无法断定李生是否为李炳的遗孤，但根据以下两点来看，这种可能性很大。首先从年龄上看，遗孤当年已经八岁，正上学念书之年龄；其次，从信中口吻看，焦循似乎觉得汪近垣资助李生，乃是理所当然之事，而近垣是李炳的门人，身份上也相符合。}
\BookNote{35.}{焦循，《李翁医记》卷上，载刘建臻点校，《焦循诗文集》(下)，第450-457页；嘉庆《重修扬州府志》卷五十四《人物·艺术》，第4219-4223页。}
\BookNote{36.}{焦循，《名医李君墓志铭》，载刘建臻点校，《焦循诗文集》(上)，第393页。}
\end{BookNotes}

% 原书第 80 页（PDF 第 93 页）
\BookPage{80}{93}
\BookSubheading{三、医技与医学思想探析}

根据现有的记载，李炳无疑是个成功的医生，目前可见的资料共留下他40个医案，其中成功治愈38例，另有两例则是成功地预测了病人的不治37。若仅就记载的内容看，比起司马迁笔下仓公成功率更高。这些案例主要都为焦循所记录，在焦循的心目，李炳医技神奇，是当时世上为数不多的医术高明者38，认为他可与清初名医汪机相比肩，而且他和汪机都年七十七而卒，「医似之，寿似之。翁固石山之后一人，翁可无憾，所惜翁之后不知何年始有翁者出耳。」39同时焦还引述当时扬州的名儒显宦江涟之口评论道:「此翁老后，不可复得，惜世人无知之者。」40不过实际上，当时的普通民众对他的医术是十分认同的，根据焦循的记载，李炳晚年多居住于邵伯，治好了相当多贫苦民众的疾病。每次去北湖看焦循，「沿湖之人邀之以小舟，相接如蚁」。41显然，江涟所谓的不为世人知之，应当是指未能在中上流社会获取声名。

以上这些论述几乎都依据于焦循等人的观察和记录，其实要真正了解一个人的医技，无疑还应该（或者说更应该）求诸其本人留下的文字。令人稍感遗憾的是，他留给后世只有一部医书《辨疫琐言》。焦循曾在其墓志铭中称他「说医之文，简而有法；间为诗歌，不甚溺也。」墓志铭中这样的说辞，我们大概可以解读为，作为一位职业的医生，他的文字简洁有度，也算不错了。只要我们通读一遍《辨疫琐言》，便不难发现，焦循的这一评论是恰当的。《辨疫琐言》的文字虽然说不上雅驯，但简洁顺畅，言之有物。虽然也不时引述《内经》和《伤寒杂病论》等经典之作，可以看出作者对这两部著作，特别是后者用力颇深，但更多的则是他自己行医治病的心得和对相关经典的思辨。作为一部医书，笔者的体会，该着最大的特点在于平直、简洁而又不乏真知灼见。从该着和其他相关记载中，我们还是可以大致看出其主要的医学思想。

\vfill
\begin{BookNotes}
\BookNote{37.}{焦循，《李翁医记》，载刘建臻点校，《焦循诗文集》(下)，第453页。}
\BookNote{38.}{焦循，《李翁医记》卷上，载刘建臻点校，《焦循诗文集》(下)，第453页。}
\BookNote{39.}{焦循，《与汪损之书》，载刘建臻点校，《焦循诗文集》（下)，第445页。}
\BookNote{40.}{焦循，《李翁医记》卷上，载刘建臻点校，《焦循诗文集》（下)，第453页。}
\BookNote{41.}{焦循，《忆书》卷五，第38页。}
\end{BookNotes}

% 原书第 81 页（PDF 第 94 页）
\BookPage{81}{94}
关于李炳的医学思想，曾有研究者就《辨疫琐言》做过专门的探讨，该论文从「《瘟疫论》理论质疑」、「《瘟疫论》用药质疑」和「遣方用药特色」等三个方面做了具体的论述42。这一研究基本是对《辨疫琐言》的具体内容一个梳理和概括，并未能从理论层次对李炳的医学思想的基本特色做出提炼。按笔者的理解，该着乃是作者依据古代经典，结合自身的医学实践，对温病学的经典名作《瘟疫论》的继承和扬弃。在总体思想上比较明显地体现出折衷、温和以及实事求是的倾向。比如在治疗的基本原则上，对于瘟疫，他既不简单认同温补之法，但也反对一味攻伐，而主张探用比较温和「清通」之法祛除温湿之邪。比如他对时行的治疫之法批评道:「彼亦明知中气亏危，正不胜邪，药则仍主达原饮、三消饮、承气等汤，不且自相矛盾耶？」43认为疫气由口鼻而入，「口气通于脾，邪从口入，必先于胃，胃者脾之表也。胃喜清通，以下行为顺」44。故主张探用「轻清以开肺舒气，芳香以醒胃辟邪」的「清通」之法。为此他特别参考经典，发明了「清气饮」这一治疫之方，「治以轻清芳香，祛浊邪而复清阳，方名清气饮。」对于该方，他「历试多人，颇有效验」45。不仅如此，大黄一般被视为攻下之药，李炳也不赞成用「寒凉峻厉」来攻下，但对并不反对使用大黄，并认为大黄为良药，原因正在于其功效，「通是也」，指出，「但用大黄，须要审人之虚实，此症属于气闭，取气以通气，每有奇验。」46在「清通」的同时，他还比较强调滋阴补胃，似乎吸收了宋元四大家中的丹溪和东垣学派的思想。他称，「疫从口入，胃经是其所舍」，十分重视对胃气的处置，反对以破气峻厉之品劫杀胃阳气，而「阳气一伤，不但变症蜂起，且恐内陷」。经典中的三承气汤乃是为阳明胃实而设，然「未必人人皆胃实症」，「今人虚者多，实者少，寒者多，热者少，温补养正而愈者，十中五六，何尝尽用承气也。疫从口人，必先于胃，…气以下行为顺，用大黄者，通其下行之道路也。」47他的这一治疗原则亦可从其他的史料中获得佐证，比如府志中他的传

\vfill
\begin{BookNotes}
\BookNote{42.}{苏颖、鞠煜洁，《李炳〈辨疫琐言〉医学思想探析》，第12-13页。}
\BookNote{43.}{李炳，《辨疫琐言》，载裘庆元辑，《珍本医书集成·内科类》，第7页。}
\BookNote{44.}{李炳，《辨疫琐言》，载裘庆元辑，《珍本医书集成·内科类》，第1页。}
\BookNote{45.}{李炳，《辨疫琐言》，载裘庆元辑，《珍本医书集成·内科类》，第2-3页。}
\BookNote{46.}{李炳，《辨疫琐言》，载裘庆元辑，《珍本医书集成·内科类》，第5页。}
\BookNote{47.}{李炳，《辨疫琐言》，载裘庆元辑，《珍本医书集成·内科类》，第8-11页。}
\end{BookNotes}

% 原书第 82 页（PDF 第 95 页）
\BookPage{82}{95}
\noindent 记中，称他好用白术，被人称为「李白术」。白术是中医中常用之药，药性苦、甘，温。归脾、胃经。具有健脾益气，燥湿利水等功效。据《医学启源》记载：「除湿益燥，和中益气，温中，去脾胃中湿，除胃热，强脾胃，进饮食，止渴，安胎。」48这一味药的特性正好与他的治疗原则颇为契合。此外，他虽然对他自己所处之方常常信心满满，一旦立定，便往往不加修改，但他对待临症处方，却极力主张以实事求是的态度，因时因地因人不同而灵活处置，切不可不问实际情况，固守成法，以一待之。比如他指出:「盖行疫之年，未必人人皆疫，亦有劳伤以及里虚里寒，伤湿、伤暑诸症夹杂其中。所谓似是而非者，全在细心体认，疫虽互相传染，医者不可为疫症所拘也。」又说:「世之宗《瘟疫论》者，十人而九。但见发热恶寒，不论何经，不论虚实寒热，春则曰春瘟，夏则曰时疫，秋则曰秋疫，冬则曰冬瘟。方则寒凉峻厉，加减出入，立案不叙脉症之理，但曰瘟疫几朝，症重防变，医者既先立不败之地，及至败坏，则案中原已载明，与医无涉，予实耻之。」49既表明了他实事求是的治疗原则，也道出了他耿直的个性。

焦循是有深厚学养的大儒，不用说是有智慧的人，而他对李炳的发自内心的信任和认同50，虽然不无两人在旨趣爱好具有类同之处（比如都对「易学」感兴趣）等方面的因素，但主要的无疑是他经过长期观察和检验逐步自然形成的。而且从李炳的著述来看，他不仅对经典著作颇有钻研，而且也有丰富的实践经验，而且更为重要的是他还能够将经典论述与自身的实践和思考很好结合起来，形成一套自己独到而平实在理的治疗原则和治疗方案。而且从当时对他多「忌之、谤之者」，也可以从另一个侧面看到他的实力和在业界的影响。因此，李炳医术的高明当无疑义。但与此同时，不可否认的是，他不仅没有闻达于当时，死后的声名，虽然由于焦循的大量记载而相对较高，但与焦循心目中可以比肩的汪机相比，仍然有极大的差距。那么其中的原由是什么呢？这大概就需要我们从李炳本身的行为和思想特色以及当时医生医名的形成和流传机制来加以探析了。

\vfill
\begin{BookNotes}
\BookNote{48.}{参阅南京中医药大学编着，《中药大辞典》（第2版），上海科学技术出版社2006年版，第942-945页。}
\BookNote{49.}{李炳，《辨疫琐言》，载裘庆元辑，《珍本医书集成·内科类》，第4、12页。}
\BookNote{50.}{焦循在李炳去世时候后，在探究易学时，有个发现与李炳从医学中得出的认识不谋而合，还甚为感叹，「惜不能起而相证也」。（焦循，《易余籥录》卷十三，载刘建臻点校，《焦循诗文集》(下)，第823-824页）焦对李挂念可见一斑。}
\end{BookNotes}

% 原书第 83 页（PDF 第 96 页）
\BookPage{83}{96}
\BookSubheading{四、为人与行医之特色}

李炳虽然医术高明，但其一生，称之为穷困潦倒，似不为过。焦循在他的墓志铭中称:「君为贫人贱士治疾，必竭尽心力，寒暑暮夜，闻召即行，而短于伺候富室显者，故身后无余财。」并在最后的铭文中言：「财利所在，让之他人。拙于求富，巧于济贫。」墓志铭中出现这样的话语，他家庭的经济状况不佳便不难想见。所以，他死后，焦循就不得不为他遗孤的抚养问题操心了51。实际上，从留下的医案中，也无意中透露出了他的贫困，比如周小濂因牙龈溃烂久治不愈，延请李炳诊治，李炳「适衣破衣」，为周所轻视52。这显然不是一个步入主流社会「名医」应有的状态，那些深受中上流社会所信任的「名医」经济状况的小康乃至富足，从明清的小说等文献中，就很容易观察到53。就是和李炳同时代的吴医顾雨（玉）田，「扬州人以千金求其一至为幸」54。对于李炳的困窘，焦循的解释是，他「短于伺候富贵显者」，「治富贵者之病，每不及治贫贱者之病之用力」55。认为其治病，不计财利，「及其愈也，所报或无一钱，君以为快」56。也就是说，李炳是个品性高洁的「喻于义」而不是「喻于利」的君子，他的贫困是因为他不屑于与富贵显者结交。这样的说法虽然不能说全是溢美之词，但作为一个职业医生，同样是给人看病，想必没有人会不愿意获取相对较高的报酬。实际上我们从现有的李炳的医案来看，他诊视的物件大多应属于有一定地位和资产之人。而且若有达官贵人礼遇信任他，他也会竭尽心力予以救治。比如道台和腾额病「两足瘠弱，不能行，以礼延翁。翁感其知己，为留三月，治之而愈」57。而且根据焦循早年的日记的记载，他不仅因为洪姓商人许以五百金的年薪而远赴楚地两年，而且还曾在给谋富人治疗久

\vfill
\begin{BookNotes}
\BookNote{51.}{焦循，《与汪损之书》，载刘建臻点校，《焦循诗文集》（下)，第445页。}
\BookNote{52.}{焦循，《李翁医记》卷下，载刘建臻点校，《焦循诗文集》(下)，第454页。}
\BookNote{53.}{这从明清时期的两部小说《壶中天》和《医界现形记》中可以得到很具体的体现。相关的内容可以参阅这两部小说以及张华最近的硕士学位论文。(《清代医家的行医之道：以〈壶中天〉与〈医界现形记〉为中心》，南开大学硕士学位论文，2010年5月）}
\BookNote{54.}{李斗，《扬州画舫录》卷十《虹桥录上》，第245页。}
\BookNote{55.}{焦循，《与汪损之书》，载刘建臻点校，《焦循诗文集》（下)，第445页。}
\BookNote{56.}{焦循，《名医李君墓志铭》，载刘建臻点校，《焦循诗文集》（上)，第395页。}
\BookNote{57.}{焦循，《李翁医记》卷下，载刘建臻点校，《焦循诗文集》(下)，第456页。}
\end{BookNotes}

% 原书第 84 页（PDF 第 97 页）
\BookPage{84}{97}
\noindent 治不愈之「右肋刺痛」病时，主动索要五百金58。焦循的解读或可以说寄予他作为一个「儒者」的美好愿望，未必全然符合李炳的实际。李炳并不见得不愿意为富贵显者诊视而获得较高的报酬，而是他不善于与他们打交道，而且也不太为中上流社会所接受和信任。而之所以造成这种状况，不在于他的医术，而是他的性情和行医方式等因素的使然。

从留存至今有关李炳的史料中，我们大概看出，他首先是一个缺乏心机、敢于任事的性情中人，同时也是一位性情孤傲、不善逢迎与变通的耿直古板之人。用今天时髦的话来说，他大概是个情商不高的人，或者说是个智商远高于情商的人。焦循曾在其墓志铭对其行医处事风格叙述道：

\begin{BookQuote}
胸有定见，不善随众浮沉，病已则戒勿药，不屑以调理为名，奔走射利。或制一方，令服百剂、数十剂不更，增损均与世俗医相反，而识者遂希。至于生死在呼吸之际，人攻君补，人塞君通，人寒君热，以口舌争之而不足，以身名性命誓而决之，手调其药而坐验其啜，不效不已59。
\end{BookQuote}

这一描述，很容易在留存至今的他的医案中找到佐证。若细心阅读他的医疗记录，往往会给人以下两点深刻的印象：一是作为医生，他缺乏自我保护意识，往往独排众议，坚持己见，为了实施自己的治疗方案不惜冒险，故被时人称之为「李大胆」60。比如，在治疗江春族人江心培的例子中，当时江感患伤寒，时医治以寒剂，不效，但李炳独自主张用温中之法，探用附子、人参等药，众医难之。李「谓诸医曰:『诸君敢包医否？』众曰:『不能。』李曰:『吾能之，吾从此居此不去，果死，治以庸医杀人罪，何如？』诸医唯唯退。乃增为人参四两，附子二两。」61在诊治周小濂的牙龈溃烂之病时，则「请屏诸医，吾独任其治，不愈，甘受其罚」62。虽然我们在现今的医案中，看到的都是成功的案例，但实际上，诊治疾病，特别是一些危急或疑难之症，不要说在当时的医疗条件下，就是当下，恐怕

\vfill
\begin{BookNotes}
\BookNote{58.}{焦循，《理堂日记》，嘉庆元年七月初十日。这些内容均不见于焦循后来撰写的《李翁医记》。}
\BookNote{59.}{焦循，《名医李君墓志铭》，载刘建臻点校，《焦循诗文集》（上)，第394页。}
\BookNote{60.}{焦循，《理堂日记》，嘉庆元年七月初十日。}
\BookNote{61.}{焦循，《理堂日记》，嘉庆元年七月初十日；李斗，《扬州画舫录》卷十二《桥东录》，第279页。}
\BookNote{62.}{焦循，《李翁医记》卷下，载刘建臻点校，《焦循诗文集》(下)，第454页。}
\end{BookNotes}

% 原书第 85 页（PDF 第 98 页）
\BookPage{85}{98}
\noindent 也不会有多少高明者，能有十足的把握。不难想见，李炳这样的做法一定也遭遇过不少失败例子。所以这样的敢于任事，在时人看来，也许就不免有些不够审慎甚至有些孟浪了。对于更加珍惜性命、也有更多选择机会的显贵们来说，自然也就会对他的治疗多一份疑惧。二是李炳不是那种善于追逐名利的人，也不是能够逢迎人的人，他往往凭着自己心情和心理感受支配自己的医疗行为，用焦循的话来说，就是「其拯人之急不畏劳烦，不恤人言，尤当于道谊学问中求之」63。比如，他直到生命的最后时候，才有幸得到子嗣，所以对于影响某人血胤传承者的疾病诊治就特别尽心，像前面谈到他对焦循之子的廷琥不计利害、弹精竭虑救治，除了有焦循多年的知遇之情这一因素外，也在于廷琥之安危，对焦循的血脉相传「所关甚重」。而且就是对于素昧平生之人，若关系血脉相继，他也会不遗余力地加以救治。乾隆五十六年（1791）冬，有一市井贫民之子得喉病，「喘促将死」，他得悉该小孩父母只有这一个儿子，便顿生恻隐之心，于是「反复求之」，最终得以获救4。以他这样的性格，和那些富贵显者交往，很难不发生矛盾，一介没有任何功名的医生，很难让大多数显贵们对他有真正的尊重，而且对自己的疾病和治疗也往往隐约有些见解，难免对他独特而又固执的意见产生想法和不放心，而另一方面若不能让他感受到尊重和信任，则又很难激发他治疗的热情。相反，为那些普通的民众治病，尽管受益微薄，但他们的尊崇和顺从则往往能带给他非常良好的心理感觉。

此外，他不知变通，「或制一方，令服百剂、数十剂不更」的做法，也有违医学处方通行的做法，自然也会让多少懂些医学的中上层社会的人士难以适用。不仅如此，由于他往往坚持己见，性情孤傲，不善交际65，也必然会使得他在同行中树敌不少，应该正因如此，焦循才会在他挽救了廷琥的生命后，感叹道:「余于此始恍然于忌之谤之者，真为庸医。」实际上，由于焦的族中的医者，每每向他批评

\vfill
\begin{BookNotes}
\BookNote{63.}{焦循，《李翁医记》卷上，载刘建臻点校，《焦循诗文集》(下)，第453页。}
\BookNote{64.}{焦循，《李翁医记》卷下，载刘建臻点校，《焦循诗文集》(下)，第457页。}
\BookNote{65.}{焦循在墓志铭中说他，「尝往来吴越、荆楚之间，所交落落」。（载刘建臻点校，《焦循诗文集》(上)，第395页）在他死后给汪损之的信中亦言:「翁平日相好无多，惟弟与兄数人。」（载刘建臻点校，《焦循诗文集》(下)，第445页)}
\BookNote{66.}{焦循，《李翁医记》卷上，载刘建臻点校，《焦循诗文集》(下)，第452页。}
\end{BookNotes}

% 原书第 86 页（PDF 第 99 页）
\BookPage{86}{99}
\noindent 李炳，也一度让焦「惑于其言」67。这些必然会使得他在当时的医界，缺乏良好的人脉关系。这样，他的不为当时的主流社会所认可和接纳，也就在所难免了。

五历史记忆

李炳的死，让焦循一家哀痛不已，在他当时写给朋友汪损之的信中称，「弟之妻女儿媳无不陨涕恸哭者」。作为一个著名的文人学者，他深知，对于一个逝者来说，欲使其生命获得延续，除了子嗣而外，最重要就是著述和事迹的流传了。所以他很快就开始着手李炳著作的收集。他跟汪损之说，《辨疫琐言》他已经整理好了，另外两部书，则请汪找出来。并「徐谋付梓」68。我们并不清楚汪损之最终是否将李炳存放他那儿的那两部著作找出来，但从此后并未见有人引用过这两部著作，且也早已亡佚这一结果看，情形恐怕不会乐观。焦循虽然学问深厚，但自己也只是一介贫士，他谋求刊印李炳著作的愿望，显然是落空了。但李炳确实一直活在他的记忆中，不仅在后来的研究和生活中会不时地想起他69，而且还用自己的笔墨为他构建一份「历史记忆」。在李翁去世前，他就撰成了《李翁医记》，其死后，则又撰写了颇具真情的「墓志铭」。同时，利用他参与府志纂修的机会，为他在府志中留下了篇幅颇巨的传记。后来，他通过医生杨赞令听到另一名医生朱培五所说有关李炳的讯息，也记录下来，编入了他的笔记《忆书》。其中云：

\begin{BookQuote}
李西垣精于医，不善逢迎，故不行于城，常往来于淮，而往邵伯镇者最久。…乙丑秋殁后，邵伯人集其十数年来所用方药，各依脉案，分门类居之成簿。有疾则设位祷于翁，检簿中方案相合者服之，顿愈。愈则谢以纸钱。70
\end{BookQuote}

虽然焦循为李炳出版著作愿望未能实现，不过他为李炳写下的这些文字有幸先后都被刊行71。而且由于焦循的声名，他的墓志铭还被多部十九世纪编纂的传记资料收

\vfill
\begin{BookNotes}
\BookNote{67.}{焦循，《李翁医记》卷上，载刘建臻点校，《焦循诗文集》(下)，第452页。}
\BookNote{68.}{焦循，《与汪损之书》，载刘建臻点校，《焦循诗文集》(下)，第445-446页。}
\BookNote{69.}{可参阅焦循，《理堂札记·丁卯札记·示季蕃》，载刘建臻点校，《焦循诗文集》(下)第655页；焦循，《易余籥录》卷十三，载刘建臻点校，《焦循诗文集》(下)，第823-824页。}
\BookNote{70.}{焦循，《忆书》卷五，第38页。}
\BookNote{71.}{参阅本文第一部分。}
\end{BookNotes}

% 原书第 87 页（PDF 第 100 页）
\BookPage{87}{100}
\noindent 入72。另外焦循最后那段记载也说明，李炳死后，他的作为医生的技艺至少在一段时间内，依然保留在邵伯乡民的记忆中。此外，焦循之子廷琥也在其著作和《辨疫琐言》的附记中多次忆及李炳，甚是感念他的医技和救命之恩73。除了焦氏父子，同邑名儒刘毓崧也在李炳去世五十年后，在已给他学生汪近垣的著作写序时，提到了他，称他为「同邑名医」，并对西垣先生的医学成绩颇有赞赏。刘毓崧生于嘉庆二十三年（1818），当时李炳已经去世八年，故他对李不可能有实际的接触，从序文中看，他对李炳的了解完全来自于焦循父子的记录以及府志中的传记74。不仅如此，应该正是缘于焦循为其所撰的墓志铭的广泛影响，刘锦藻完成于民初的《清续文献通考》还在「经籍考」中著录了他的《金匮要略注》75，尽管该书早已亡佚。

显而易见，我们今天能了解李炳的基本讯息，几乎要完全归功于焦循的记录和对其《辨疫琐言》的整理和誊抄，若不是他能在51岁那样有幸结识焦循，这位医技高明却不为主流社会认可的医生恐怕早已湮没无闻。不过，尽管焦循学术影响广泛而深远，他的著作也不缺乏读者，但有关李炳那些记录在其博大巨集富的著述中，也不过是并不起眼的一小部分，若不是对医学所有兴趣者，甚至可能不会留下什么印象。因此，作为一位并无其他特别才华的职业医生，如果未能留下自己的医学作品，即便后人对焦循再关注，李炳至多也只不过是一个缺乏主体性的附庸而已。不过有幸，他不仅写有著作，而且还留有弟子。虽然周金声在他死后也很快去世，但汪近垣在在他死后，仍行医五十余年，且颇有声名。从刘毓崧的说法来看，汪不仅深得李炳真传，而且也在其著作中极力阐扬其师的经义。因此，可以说他的医学生命，至少在一段时间内获得了一定的延续。不过，尽管汪生前的境遇似乎要明显好于其师，但他显然也不是能光大其师学说的名医，在现存的文献中，就管见所及，

\vfill
\begin{BookNotes}
\BookNote{72.}{如收录于钱仪吉编纂、成书于道光年间的《碑传集》（卷147，中华书局1993年标点本，第12册，第4347-4349页）和李恒编纂、成书于光绪十一年（1883年）的《国朝耆献类征初编》（卷483，江苏广陵古籍刻印社1990年版，第20册，第335-336页）}
\BookNote{73.}{焦廷琥，《蜜梅花馆文录·郑素圃医案序》，丛书集成初编本，第16页。李炳，《辨疫琐言》，「焦廷琥附识」，载裘庆元辑，《珍本医书集成·内科类》，第12页。}
\BookNote{74.}{刘毓崧，《通义堂文集》卷十一《汪近垣先生金匮要略阐义序》，第512-513页。}
\BookNote{75.}{刘锦藻，《清续文献通考》卷275《经籍考十九》。从刘氏对该书的按语可以看出，他所以著录该书，应是对其所谓的「肝之本在右边而行于左」之说颇有感触，不过从行文看，则基本转录自焦循所撰的墓志铭。}
\end{BookNotes}

% 原书第 88 页（PDF 第 101 页）
\BookPage{88}{101}
\noindent 除了刘毓崧的序言，并未发现再有其他有关他的记载，我们甚至都不知道他的名字76。虽然他也在其师论着的基础上写成了《金匮要略阐义》，而且其子也积极准备刊刻出版，但结果也未能如愿。至今仅留下一部抄本7。

李炳死后，虽然在十九世纪，主要因为焦循的缘故，在当时的一些文献中留下了李炳的足迹，但随着时间推移，由于他以及弟子汪近垣的作品均未能刊刻出版，拥有高明医技和独到医学思想的医生李炳，也就渐渐不再为世人熟悉和记忆了。检索道光以后的历史文献，有关李炳的记载，不过仅有的几条，而众多的医学著作，也几未见有对李炳的引用和评论。1929年，江甯的汪绍达先生在上海编辑刊刻了「回溯社医书四种」，《李翁医记》即为其中的一种，该书在重版时，编者在正文前加了一个序言。虽然汪也算是李的同乡，但却完全不知道李翁是何人。他在序言中首先就说:「李翁，不知何许人，要为乾嘉时淮扬一带之高手也。惜其名字不传，无由识其生平大概。」不过数年后，李炳的命运似乎又一次发生转变。1936年，近代名医、藏书家、出版家绍兴的裘庆元先生从他所藏的众多医书中，选取较实用的精本、孤本、抄本、未刊稿等九十种加以分门别类，编纂成「珍本医学丛书」，由上海世界书局出版。李炳的《辨疫琐言》有幸作为12中「内科类」著作中的一种被刊出。该书对每一种著作均写有一提要，编纂者在写该着的提要时，显然参考了府志中的传记和焦循所撰的墓志铭，准确地介绍了李炳的基本生平。对于著作，提要评论道:「撰者感于世人治疫，皆为吴又可《瘟疫论》所惑，对吴氏论疫提出异议。…用意周微，亦具至理，洵为辨疫名言。」79由于裘庆元在当时中国中医学界的重要影响，以及该丛书所收之书均是编者认为是医书中流传少而内容精的珍品，因此《辨疫琐言》乃至李炳自然就开始日渐受到关注了。特别是随着温病学的发展愈益被视为明清医学发展的最重要的成就之一，李炳及其《辨疫琐言》也就更是不断被人提及了。一些温病学和医史论着，往往将李炳的《辨疫琐言》当作

\vfill
\begin{BookNotes}
\BookNote{76.}{从刘毓崧的序言看，近垣应是他的号。他的这一号或许跟其师的号西垣有关。}
\BookNote{77.}{现存南京图书馆。}
\BookNote{78.}{汪绍达，《李翁医记叙》，1929年回溯社刻本，第1a页。}
\BookNote{79.}{裘庆元辑，《珍本医学丛书·内科类》之《书目提要》，第2-3页。}
\end{BookNotes}

% 原书第 89 页（PDF 第 102 页）
\BookPage{89}{102}
\noindent 吴有性《瘟疫论》后温病学领域颇具影响的代表性著作8°。不过就像我们已经看到的，李炳《辨疫琐言》的重要影响并没有发生李炳生活的时代以及他身后相当长的时间里，而是随着该着被纳入「珍本医书集成」出版后，出现在现代温病学和医史学研究者的认识中。

如今，李炳已是在中国医学史特别是温病学史占有一席之地的名医，业已成为中国医史和温病学研习者比较需要面对的记忆物件，并且由于焦循称其「为贫人贱士治疾，必竭尽心力，寒暑暮夜，闻召即行，而短于伺候富室显者」，所以偶尔也会作为具有高尚医德的名医被记忆81。与此同时，焦循乃中国学术思想史和地方史研究中的重要人物，随着学界对焦循研究展开，李炳也往往因其与焦循的人生密切相关而被论及82。李炳身后的历史记忆，自然不可能脱离他生前的作为，但所有这些记忆，显然都与焦循的记录以及《辨疫琐言》有幸被纳入「珍本医书集成」出版有关，若没有这两个不无一定偶然性的条件，一代「名医」恐怕早已湮没在历史的长河中了。

\vfill
\begin{BookNotes}
\BookNote{80.}{比如一部颇具权威性的温病学著作写道:「《温疫论》问世后，研究温疫的医家和著作相继涌现，其中影响较大的如戴天章《广温疫论》、余师愚《疫疹一得》、刘松峰《说疫》、熊立品《治疫全书》、陈耕道《疫烦草》、李炳《辨疫琐言》、汪期莲《温疫汇编》等。」（杨进主编，《新编温病学》（上册），学苑出版社2003年版，第225页）这样的说法在现代中医学或医史学著作中，时可见到。（比如邱模炎等主编，《中医疫病学》，中国中医药出版社2004年版，第14页；刘镜如编着，《中医史话》，甘肃人出版社1981年版，第110页，等等)}
\BookNote{81.}{参见孙溥泉、李恩昌、孙健慧，《清代名医的医德（一)》，《中国医学伦理学》2008年第6期，第150转158页。}
\BookNote{82.}{可参阅最近两部有关焦循的著作，刘瑾辉，《焦循评传》，广陵书社2005年版，第29、197、201页；刘建臻，《焦循著述新证》，社会科学文献出版社2005年版，第29、69-75、308页。不过这两部论着中的相关论述都不无可以商榷之处，比如前者第29页称焦循「在《北湖小志》中专门作《李翁医记》二卷」，有误，《李翁医记》完成于嘉庆十年，早于《北湖小志》，实际上是两书同时刊刻出版。而且作者似乎并不清楚李翁即李炳。第201页，将汪绍达于1929年（原序题己巳年）作《李翁医记叙》一事归入嘉庆十四年（1809年）之目下。后者附录的《焦循著述系年》，在乾隆五十五年（1790年）条目下，记：「十月，病呕血且咳，为李翁治愈，信其医术之神。」(第308页）但实际上，《李翁医记》明确说，那次是他第二次请李翁看病，当时「犹未信翁之神也。真正完全信服，乃十五年之后的事。}
\end{BookNotes}

% 原书第 90 页（PDF 第 103 页）
\BookPage{90}{103}
六结语

我们今天称李炳为清代扬州「名医」，大概不会有任何人会表示异议，然后，从前面的论述不难看到，若回到李炳生活的时代，恐怕很多人未必认同。尽管焦循盛赞其医术，但他也无法讳言，李并不见重于当世，即不为主流社会所认可和信任，贫困潦倒，全然没有当世名医的气派。李炳的失意，不在于他的医术，而是他过于耿直和自信、不会逢迎和讨巧以及不善交际的性格的使然。在清代，国家并无医生的准入制度和考察医生水准的评判标准，医生的是否能够立足，几乎完全取决于自己能否在医疗市场占有一席之地。对于医生的成名来说，除了一定的医技外，同样重要的还有其似乎擅长自我宣传炒作、是否善于揣摩病人及其家属心思和逢迎讨巧83。此外，医生的文学修养及其与士人的唱和对其声名的抬升也具有重要的意义，而若能让那些有名的文人为自己写下一些文字，让无疑更有利于为扩大自己的影响加分。笔者在翻阅明清的文集时，常常可以看到那些著名文人为一些医生写的各种序文，写作的缘由或为受朋友的委托（也就是说某些医生托人相求），或者是因为其为该文人本人、亲友治好了疾病，坚持不要报酬，而只是希望文人为自己写一篇序文。这些医生显然是希望通过与文人的交往，和文人为他们写序文，让自己变得更有文化，更易被视为受主流社会认可的「儒医」以及更容易让自己的声明传之后世。但这样一些追求和能力，除了医技，李炳似乎都不具备。一方面，他既不会也不愿去宣传和炒作自己，以及去琢磨和讨好病人及其家属；另一方面，他虽然识文断字，但一生并无获取任何功名，文字能力似乎也不突出，而且似乎也不愿刻意去跟那些文人学士交往，也不会希图借助著名文人的文字来扩展自己的影响。这从他跟焦循的交往中不难看出，虽然随着他们交往的增加，焦循的学术声望不断提高，他们感情也不断深入，但李炳却好像从来没有让焦循为其写什么文字，连自己著作的序文都没让写。而焦循为他留下的那些文字，也完全是出于自己内心情意，除了《李翁医记》是在他临去世之前所写外，其余都是他去世之后才撰着的。李炳能够遇到焦循，或许可以说是个意外，而他最终因为焦循的文字以及帮他整理著作而在历史上留下声名，大概也算是「无心插柳柳成荫」吧。另外，《辨疫琐言》的抄本能够为裘庆元收藏，并最终入选《珍本医书集成》，似乎也是偶然的幸运，比

\vfill
\begin{BookNotes}
\BookNote{83.}{参阅张华，《清代医家的行医之道:以〈壶中天〉与〈医界现形记〉为中心》，第15-80页。}
\end{BookNotes}

% 原书第 91 页（PDF 第 104 页）
\BookPage{91}{104}
\noindent 起他生前的境遇，他身后的遭遇显然要幸运得多。但不管他是如何幸运，若没有生前高明的医技做基础，那其他的一切也就无从谈起了。

嘉庆十年（1805）七月，医技高明的医生李炳走完了贫困了的一生，却幸运地在身后获得了较高的声誉。如今，二百多年过去了，在历史的记忆中，他是著名学者焦循的朋友、恩人，也是具有一定建树的温病学家，还是医德高尚的「名医」。所有这些，其实都不过是后人从自己意愿和角度出发，为他贴上的标签。不管怎样，他首先是一个有自己思想和性情的人，无论我们有怎样的意愿和立场，我们都有必要尊重这一点。故此，笔者撰成此文，希望以李炳为中心，从社会文化史的角度，在有关他的历史记忆中绘上新的一笔。


% 原书第 92 页（PDF 第 105 页）
\BookPage{92}{105}
